% Section: HygieneBench Design
% Source: paper_draft_v0.1.md § 3. HygieneBench Design
% [VERIFY ACM TEMPLATE BEFORE SUBMISSION]

\section{\hygienebench Design}
\label{sec:design}

\subsection{Schema and Synthetic Generator}

\hygienebench v0.1 uses a schema with eleven entity/event tables and one anomaly
label table (Table~\ref{tab:schema}). The generator (\texttt{SyntheticHygieneGenerator})
is implemented in Python using NumPy and Pandas, with all randomness gated through
\texttt{numpy.random.default\_rng(seed)}. Entity identifiers are deterministic UUIDs
derived from \texttt{hashlib.md5(seed\_tag + index)}, ensuring reproducibility
across platforms.

% Table 1: HygieneBench schema
% Source: paper_draft_v0.1.md Table 1
% [VERIFY ACM TEMPLATE BEFORE SUBMISSION]

\begin{table}[t]
\centering
\caption{\hygienebench v0.1 schema (medium scale, $n$=1,000 users).}
\label{tab:schema}
\small
\begin{tabular}{@{}lr>{\raggedright\arraybackslash}p{3.4cm}@{}}
\toprule
Table & Rows & Description \\
\midrule
\texttt{users}                     & 1,000      & AD user accounts with identity state \\
\texttt{groups}                    & 65         & AD security/distribution groups \\
\texttt{computers}                 & 830        & Managed endpoints \\
\texttt{assets}                    & 830        & Asset inventory entries \\
\texttt{login\_events}             & $\sim$249k & Auth.\ events (30-day window) \\
\texttt{group\_membership\_events} & $\sim$51   & Group add/remove events \\
\texttt{account\_lifecycle\_events}& varies     & Account create/disable/enable \\
\texttt{endpoint\_patch\_state}    & 830        & Per-endpoint patch compliance \\
\texttt{vulnerability\_records}    & $\sim$3.5k & Per-endpoint CVE exposure \\
\texttt{remediation\_events}       & varies     & Patch and remediation actions \\
\texttt{telemetry\_freshness\_log} & varies     & Per-entity data-source freshness \\
\texttt{anomaly\_labels}           & 110        & Ground-truth labels (task+split) \\
\bottomrule
\end{tabular}
\end{table}


\paragraph{Structural priors.}
Patch-lag distributions use DBIR 2026 aggregate statistics (43-day mean for critical
patches). CVE severity distributions use NVD base score statistics. Telemetry
refresh cadence requirements follow CISA BOD 23-01 (14-day asset discovery, 72-hour
patch data)~\cite{cisa_bod2301_2022}. AD group-size distributions lack a
peer-reviewed published source and are modeled via disclosed assumptions (power-law
approximation) with sensitivity sweeps. All structural priors are documented in the
dataset card.

\paragraph{Conditions.}
Five telemetry conditions parameterize freshness and missingness:

\begin{itemize}[leftmargin=*]
  \item \textbf{C-BASE:} No artificial staleness; all sources current.
  \item \textbf{C-FRESH:} Elevated freshness flagging; sources verified current.
  \item \textbf{C-STALE:} 20\% of entities have heavy-stale telemetry ($>$14-day gap).
  \item \textbf{C-MISS:} One data source absent for 15\% of entities.
  \item \textbf{C-UNSUP:} C-BASE with anomaly labels withheld from the training pipeline.
\end{itemize}

Each condition is generated for three seeds (42, 137, 2024), producing 15 dataset
instances (5 conditions $\times$ 3 seeds) with $n$=1,000 users per instance.

\subsection{Anomaly Taxonomy}

\hygienebench defines twelve anomaly classes (Table~\ref{tab:taxonomy}), each
mapped to an ATT\&CK enabling condition (not a detection claim for any specific
technique). Enabling-condition mappings are annotated with a required disclaimer:
\emph{correlation with an ATT\&CK enabling condition does not imply that any
specific attack technique was executed or detected.}

% Table 2: Anomaly taxonomy
% Source: paper_draft_v0.1.md Table 2
% [VERIFY ACM TEMPLATE BEFORE SUBMISSION]

\begin{table}[t]
\centering
\caption{HygieneBench anomaly classes. ATT\&CK mappings are enabling-condition
  correlations; they do not imply detection of any specific technique.}
\label{tab:taxonomy}
\small
\begin{tabular}{@{}llll@{}}
\toprule
Class & Code & ATT\&CK enabling condition & Tasks \\
\midrule
Stale privileged account       & AH-01 & Valid Accounts: Domain (T1078.002)    & T1 \\
Privilege escalation drift     & AH-02 & Valid Accounts, Domain Policy Mod.    & T7 \\
Group membership drift         & AH-03 & Domain Policy Modification (T1484)    & T2 \\
Dormant account reactivation   & AH-04 & Valid Accounts: Domain (T1078.002)    & T6 \\
Impossible/unusual login       & AH-05 & Valid Accounts, Lateral Movement      & T1,T3 \\
Endpoint--identity risk corr.  & AH-06 & Exploit Public-Facing App             & T3 \\
Patch noncompliance cluster    & AH-07 & Exploit Public-Facing App (T1190)     & T5 \\
KEV exposure aging             & AH-08 & Exploit Public-Facing App (T1190)     & T5 \\
Asset inventory mismatch       & AH-09 & Defense Evasion enabling              & T4 \\
Missing endpoint agent         & AH-10 & Defense Evasion enabling              & T4 \\
Telemetry missingness cluster  & AH-11 & Defense Evasion enabling              & T4 \\
Abnormal remediation delay     & AH-12 & Exploit Public-Facing App             & T5 \\
\bottomrule
\end{tabular}
\end{table}


\subsection{Benchmark Tasks}

Seven tasks cover the twelve anomaly classes across different entity scopes and
temporal windows (Table~\ref{tab:tasks}). Task injection counts are per dataset
instance; injection uses \texttt{AnomalyInjector(seed + 999)} to ensure a distinct
RNG stream from the generator.

% Table 3: Benchmark tasks
% Source: paper_draft_v0.1.md Table 3
% [VERIFY ACM TEMPLATE BEFORE SUBMISSION]

\begin{table*}[t]
\centering
\caption{Benchmark tasks. $k$ values are pre-registered per TASK\_SPECS.md and
  reflect operational SOC review budgets.}
\label{tab:tasks}
\small
\begin{tabular}{@{}lllrrl@{}}
\toprule
Task & Description & Entity scope & $k$ & Imbalance & Anomaly classes \\
\midrule
T1 & Stale privileged account risk          & Privileged users           & 10 & 1:50  & AH-01 \\
T2 & Group membership drift                 & Users with group events    & 20 & 1:167 & AH-03 \\
T3 & Endpoint--identity risk correlation    & Users with endpoints       & 15 & 1:333 & AH-05, AH-06 \\
T4 & Telemetry coverage gaps                & All computers/assets       & 20 & 1:14  & AH-09, AH-10, AH-11 \\
T5 & Patch--vulnerability hygiene           & All computers              & 25 & 1:59  & AH-07, AH-08, AH-12 \\
T6 & Dormant account reactivation           & Users with lifecycle events & 10 & 1:250 & AH-04 \\
T7 & Escalation drift detection             & Users with group add events & 10 & 1:500 & AH-02 \\
\bottomrule
\end{tabular}
\end{table*}


\paragraph{Dataset splits.}
Labels are stratified 60/20/20 by anomaly class. Anomalous entities are assigned to
the test split, with normal entities split randomly across train/val/test. The split
is deterministic given the seed.
